\section{Cognitive-Inspired Tasks}
\label{sec:tasks}

To demonstrate the usefulness of \TVD in the development of algorithms for real-world applications, we conducted experiments on three representative tasks. Specially, we argue that they are closely related to the development milestones of human intelligence. To be exact, at the age of 3, children start to occupy the capability of discriminating the genres of TV content~\citep{wright1994young}. By 10 years old, as memory grows, children can follow major plot ingredients presented explicitly, thus being able to identify one cinema/TV modality (e.g., video) given another (e.g., text)~\citep{gunter2005children}. Writing empathic and context-aware plot descriptions, however, is a complicated skill which requires not only intelligence but also proper training, so most children have trouble mastering it until the end of the Formal Operational Stage (around 16 years old)~\citep{piaget2003psychology}. Analogously, by visiting the tasks of Genre Classification, Plot Retrieval, and Plot Text Generation, one can also study multimodal systems from a shallow level to a deep level, thus revealing the weaknesses and limitations of tested algorithms.

\noindent\textbf{Genre Classification.} Genre Classification is the task of labeling a sample with one or more genre tags, which is similar to the multi-classification problems concerned by previous works in the cinema domain~\citep{moviescope2019arxiv,movienet2020eccv}.
Since the plot already provides concentrated genre information (e.g., fighting plots in action dramas, dating plots in romance dramas), we take the plot-oriented clips as the primary unit to be processed.
As for the metric, following common practice in multi-classification, we calculated the F1 scores. Besides the Micro setup, considering the unbalanced genre distribution of clips (see Fig.~\ref{fig:genre_tags}), we also benchmark the Macro one.


\noindent\textbf{Plot Retrieval.} This task asks a model to retrieve matched results whose modality is different from that of the input~\citep{cmr2016arxiv}.
As all modalities of \TVD are made parallel, it serves as a good test bed for cross-modal retrieval.
Specifically, in this work, we experimented with two different retrieval settings: Image-Text, and Video-Text. Among them, the video mode contains both visual signals and text signals (i.e., image frames and subtitle texts).
For the text side, we used plot-focused sentences rather than the BSCs. The main reason is that a BSC can be quite general, e.g., ``\textit{Actor A is so beautiful!}'', in which case it may point to a huge number of correct video/image candidates, making the evaluation difficult. 
Additionally, we did not consider the retrieval between image and video as the former is directly captured from the latter.

We reported two groups of results for the Plot Retrieval task. The first is Hit@$n$ ($n \in \{1,5,10\}$), i.e., the fraction of correct answers that rank in the top $n$ retrieved samples.
The second is Matching Accuracy (following \citep{meter2022cvpr}). To be exact, we asked the model to traverse all possible pairings between samples from two modalities and calculated the overall accuracy.

\noindent\textbf{Plot Text Generation.} Beyond the above two understanding-only tasks, we added this more complex generation problem.
In the cinema/TV domain, previous studies have attempted to develop methods that can take video clips (including frames and subtitles) as input and produce scene descriptions as output, so as to help visually impaired individuals enjoy movies~\citep{vc12017itm,vc22020cvpr}. Our Plot Text Generation task is more challenging than that, as the tested model is additionally expected to perceive the dramatic context and the mental states of roles. To the best of our knowledge, this is one of the most ambitious multimodal text generation goals in the community~\citep{wang2021survey,zelaszczyk2023cross}.
In our study, we selected three popular types of automatic metrics: BLEU (@$\{1,2\}$)~\citep{bleu2002acl}, METEOR~\citep{meteor2005acl}, and CIDEr~\citep{cider2015cvpr}. 

% Apart from generating plot text, for more comprehensive evaluations we trained the system to write BSCs. For each clip, there exists one single ground-truth plot-focused sentence, but there are many BSCs.
% Consequently, we regarded the top 10 upvoted ones as the golden references as they are typically informative and of high quality.


